\section{The consistency proof}
\label{sec:proof}

This section follows the argument from the evaluator up to the final theorem.
The intent is to make the shape of the proof and its load-bearing lemmas
visible; the detailed proofs are left to a later instalment.

\subsection{Shape of the argument}

Everything reduces to one implication, called the soundness bridge:

\begin{center}
if \isa{p} checks as a proof of \isa{J}, and the hypotheses of \isa{J} are
satisfied by an assignment \isa{A}, then the conclusion of \isa{J} is satisfied
by \isa{A}.
\end{center}

\noindent
Given it, \S\ref{sec:skeleton} closes the argument in half a page.  Proving it
requires four things, and they are the subsections that follow.

\begin{enumerate}
\item The step-indexed evaluator determines at most one value, so that
      \isa{evals} behaves as a partial function and satisfaction is well
      behaved.
\item Every branch of \isa{valid\_step} preserves satisfaction, assuming that
      every judgment in the tail of the proof does.
\item That local statement can be lifted along the proof list by induction.
\item Satisfaction of an encoded equation and of an encoded disequality can be
      unpacked into statements about values.
\end{enumerate}

\noindent
A remark on technique before the details.  \QGA{} induction instantiates a
schematic predicate \isa{Q}, so a statement can only be generalized before
inducting if the generalization is part of the object-level formula.  Several
proofs in this file are therefore stated internally in the form
\isa{∀u. ∀A2. …}, or as nested implications \isa{… ⟶ (… ⟶ …)}, and unpacked
afterwards with \isa{forallE} and \isa{implE}.  \isa{eval\_fuel\_N} is the first
instance: the induction on fuel runs on \isa{∀u. ∀A2. eval\_fuel k u A2 N} and
the desired statement is recovered by two instantiations.

Nesting an implication inside a motive has a price.  $\longrightarrow$ introduction
requires its antecedent to be grounded, so every antecedent of such a motive
needs a groundedness lemma.  The checkers all have one.  Satisfaction does not,
and \S\ref{sec:satat} is what supplies it.

\subsection{The evaluator is a partial function}

The fuel evaluator is total as a function into \isa{num}, which is the
statement \isa{eval\_fuel\_N: ⟦k N; t N; A N⟧ ⟹ eval\_fuel k t A N}, and $0$ is
its ``no result'' value.  Success is preserved when the budget grows.

\thy{bga-fuel-add}

\noindent
This is proved by induction on the increment, from the one-step version
\isa{eval\_fuel\_success\_suc}.  Two successful runs with different budgets
agree, since both can be raised to the common budget \isa{k + l}.

\thy{bga-fuel-unique}

\noindent
The existential form follows, and it is the assumption \isa{evals\_det} that
the abstract locale of \S\ref{sec:skeleton} demanded.

\thy{bga-evals-functional}

\noindent
The remaining facts about the evaluator are one lemma per constructor, in two
directions each: from a successful run one can read off the value of the
subterms (\isa{eval\_fuel\_suc\_value}, \isa{eval\_fuel\_pred\_value},
\isa{eval\_fuel\_ifz\_zero\_value}, \isa{eval\_fuel\_app\_value}), and from
successful runs of the subterms one can build a successful run of the whole
(\isa{eval\_fuel\_sucI} and its relatives).  Together they let \isa{evals} be
used as an inductively defined big-step relation, which is how the rule
soundness proofs consume it.

\subsection{Soundness of a single step}

The central local statement is the following.  It says that a checked step is
sound relative to the tail of the proof, and it is what the induction of
\S\ref{sec:list} will feed itself.

\thy{bga-valid-step-sound}

\noindent
Its proof is a case analysis over the branches of \isa{valid\_step}, using
\isa{cases\_bool} on each guard in turn, with one lemma per branch.  All the
branch lemmas have the same premises as above and the same conclusion; they
differ only in which checker they assume.

\begin{center}
\begin{tabular}{@{}ll@{}}
\toprule
\textbf{Branch} & \textbf{Lemma}\\
\midrule
hypothesis          & inline, from \isa{sat\_hyp\_fuel\_mem}\\
cut                 & \isa{check\_cut\_sound\_fuel}\\
substitution of equals & \isa{check\_subst\_sound\_fuel}\\
induction           & \isa{check\_ind\_sound\_fuel}\\
application         & \isa{check\_app\_sound\_fuel}\\
weakening           & \isa{check\_struct\_sound\_fuel}\\
the seven equality rules   & \isa{check\_eq\_rules\_sound\_fuel}\\
the four disequality rules & \isa{check\_neq\_rules\_sound\_fuel}\\
\bottomrule
\end{tabular}
\end{center}

\noindent
Two of these carry the real content.

\paragraph{The template lemma.}
Substitution of equals, application and induction all reduce to the same
semantic fact: if a template \isa{p} instantiated at \isa{a} is satisfied, and
\isa{a} and \isa{b} evaluate to a common value, then the same template
instantiated at \isa{b} is satisfied.

\thy{bga-template-sound}

\noindent
It rests on a pair of substitution lemmas that relate syntactic substitution to
updating the assignment.

\thy{bga-subst-down}
\thy{bga-subst-up}

\noindent
Given \isa{a} and \isa{b} that both evaluate to \isa{v}, the first pushes
satisfaction of \isa{p[i ↦ a]} under \isa{A} down to satisfaction of \isa{p}
under \isa{asn\_put A i v}, and the second lifts it back up to
\isa{p[i ↦ b]}.  The search functions are then sound because they report
success only when such a template exists, which is the content of
\isa{check\_template\_sound\_fuel} and of

\thy{bga-find-phi-sound}

\noindent
The premise \isa{prev} is the induction hypothesis of \S\ref{sec:list} in the
form the search needs: every judgment occurring in \isa{rest} is sound.  The
premise \isa{sub} records that the pointer walks a sublist of \isa{rest}, which
is what keeps the recursion inside the part of the proof already known to be
sound.  Application is the same lemma applied to the template
$\varphi[\,\mathrm{body}_d[v_0 \mapsto x, v_1 \mapsto y]\,]$, with the two
habeas quid premises supplying the values that \isa{subst\_body} needs.

\paragraph{The induction rule.}
Soundness of the induction rule is the one place where a semantic induction is
performed, over the value of the term the rule inducts on.

\thy{bga-nat-ind-sound}

\noindent
Read semantically: if the hypothesis list is satisfied, the base instance
$p[v_i \mapsto 0]$ is satisfied, and the step is satisfied for every value
\isa{m} put into the assignment at index \isa{i}, then the instance at \isa{a}
is satisfied whenever \isa{a} evaluates to \isa{n}.  The proof establishes
\isa{⋀m. m N ⟹ sat\_fuel p (asn\_put A i m)} by \QGA{} induction on \isa{m},
with the base case supplied by \isa{sat\_fuel\_subst\_FD} applied to the
encoded zero, and then instantiates it at \isa{n}.  The freshness premise
\isa{fresh\_H i G} is what
allows the assignment to be modified at \isa{i} without disturbing satisfaction
of the hypotheses.  \isa{check\_ind\_sound\_fuel} then supplies the three
premises from the corresponding judgments found in the proof list.

\subsection{From steps to proofs}
\label{sec:list}

The lifting step is a dedicated induction principle over proof lists.

\thy{bga-check-list-induct}

\noindent
It is stated for \isa{sat\_fuel} and \isa{sat\_hyp\_fuel}.  Its only caller
instantiates it at those two, and the indexed predicates it uses internally are
defined in terms of them.

Internally it is proved by \isa{list\_induct} on a statement of the form
\begin{center}
\isa{∀A. ∀K. ∀k. check\_list pf ⟶ (mem K pf ⟶ (sat\_hyp\_at k (hyp\_of K) A ⟶ sat\_fuel (conc\_of K) A))}.
\end{center}
The predicate being inducted on has to be a single formula, so the
generalizations over the assignment, over the judgment and over the budget sit
inside it as \QGA{} quantifiers.

The three antecedents are why the budget is explicit.  Each one has to be
grounded before $\longrightarrow$ can be introduced, and \isa{check\_list\_bool},
\isa{mem\_bool} and \isa{sat\_hyp\_at\_bool} supply exactly that at the three
levels.  The conclusion needs no such treatment, since $\longrightarrow$ introduction
constrains only the antecedent, which is why \isa{sat\_fuel} appears there
unindexed.  With \isa{sat\_hyp\_fuel} in the antecedent instead, the introduction
would be blocked.

The indexed predicate never leaves the proof.  The statement of the lemma is in
\isa{sat\_hyp\_fuel} and \isa{sat\_fuel}, and the bridges of \S\ref{sec:satat}
convert at the boundary.  \isa{sat\_hyp\_atD} runs downward at the two points
where the induction hands control to \isa{base} and \isa{step}, both of which are
stated in the unindexed predicates.  \isa{sat\_hyp\_at\_exI} runs upward twice,
once in the tail case where the induction hypothesis needs a budget for the
hypotheses of another judgment, and once at the end where the lemma's own
assumption has to produce one.

Instantiating the principle and discharging the step case with
\isa{valid\_step\_sound\_fuel} gives soundness for every judgment occurring in a
checked list.

\thy{bga-check-list-sound}

\noindent
Finally \isa{is\_valid\_proof} is unfolded.  Its two guards are decided with
\isa{cases\_bool}, the head of the proof is shown to be a member of it by
\isa{cons\_reconstr} and \isa{mem\_cons\_head}, and the previous lemma applies.

\thy{bga-bridge}

\subsection{Consistency}
\label{sec:consistency}

The two formula constructors are defined, and their satisfaction is unpacked.

\thy{bga-mk-eq}
\thy{bga-sat-eqE}
\thy{bga-sat-neqE}

\noindent
Both unpacking lemmas go through \isa{tag\_pack\_F} and \isa{load\_pack\_F} to
recover the two sides from the code, and then through the definition of
\isa{sat\_fuel}.  Groundedness of proof checking, the last assumption of the
skeleton, is \isa{proof\_is\_bool}, obtained from \isa{check\_list\_bool} by
list induction and from a chain of \isa{condTB} steps over the branches of
\isa{valid\_step}.

With these in place, \isa{bga\_fuller}, the union of the step-indexed semantics
with the proof checker, is declared a sublocale of \isa{consistent}.

\thy{bga-sublocale}

\noindent
Each obligation is discharged by one of the lemmas above:
\isa{mk\_eq\_N'} and \isa{mk\_neq\_N'} for habeas quid,
\isa{proof\_is\_bool} for groundedness of the checker,
\isa{sat\_hyp\_fuel\_nil} for the empty hypothesis list,
\isa{evals\_functional} for determinism,
\isa{sat\_fuel\_mk\_eqE} and \isa{sat\_fuel\_mk\_neqE} for the two unpacking
assumptions, and \isa{soundness\_bridge\_fuel} for soundness.  The theorem
\isa{syntactically\_consistent} of \S\ref{sec:skeleton} becomes available
inside \isa{bga\_fuller}, and it remains to exhibit an encoding that satisfies
the interface.
